%%
%% This is file `sample-acmsmall-conf.tex',
%% generated with the docstrip utility.
%%
%% The original source files were:
%%
%% samples.dtx  (with options: `all,proceedings,bibtex,acmsmall-conf')
%% 
%% IMPORTANT NOTICE:
%% 
%% For the copyright see the source file.
%% 
%% Any modified versions of this file must be renamed
%% with new filenames distinct from sample-acmsmall-conf.tex.
%% 
%% For distribution of the original source see the terms
%% for copying and modification in the file samples.dtx.
%% 
%% This generated file may be distributed as long as the
%% original source files, as listed above, are part of the
%% same distribution. (The sources need not necessarily be
%% in the same archive or directory.)
%%
%%
%% Commands for TeXCount
%TC:macro \cite [option:text,text]
%TC:macro \citep [option:text,text]
%TC:macro \citet [option:text,text]
%TC:envir table 0 1
%TC:envir table* 0 1
%TC:envir tabular [ignore] word
%TC:envir displaymath 0 word
%TC:envir math 0 word
%TC:envir comment 0 0
%%
%% The first command in your LaTeX source must be the \documentclass
%% command.
%%
%% For submission and review of your manuscript please change the
%% command to \documentclass[manuscript, screen, review]{acmart}.
%%
%% When submitting camera ready or to TAPS, please change the command
%% to \documentclass[sigconf]{acmart} or whichever template is required
%% for your publication.
%%
%%
%\documentclass[acmsmall,screen,anonymous,review]{acmart}

\documentclass[acmsmall,screen,review,anonymous]{acmart}

%\documentclass[acmsmall]{acmart}

%%
%% \BibTeX command to typeset BibTeX logo in the docs
\AtBeginDocument{%
  \providecommand\BibTeX{{%
    Bib\TeX}}}

%% Rights management information.  This information is sent to you
%% when you complete the rights form.  These commands have SAMPLE
%% values in them; it is your responsibility as an author to replace
%% the commands and values with those provided to you when you
%% complete the rights form.
\setcopyright{acmlicensed}
\copyrightyear{2026}
\acmYear{2026}
\acmDOI{XXXXXXX.XXXXXXX}
%% These commands are for a PROCEEDINGS abstract or paper.
\acmConference[Conference acronym 'XX]{Make sure to enter the correct
 conference title from your rights confirmation email}{June 03--05,
  2026}{Woodstock, NY}
%%
%%  Uncomment \acmBooktitle if the title of the proceedings is different
%%  from ``Proceedings of ...''!
%%
%%\acmBooktitle{Woodstock '18: ACM Symposium on Neural Gaze Detection,
%%  June 03--05, 2018, Woodstock, NY}
\acmISBN{978-1-4503-XXXX-X/2018/06}

\usepackage{enumitem}
\usepackage{lipsum}
\usepackage{array}
\usepackage{wrapfig}
\usepackage{subcaption}
\usepackage{pgfplots}
\usepgfplotslibrary{groupplots}
\pgfplotsset{compat=1.17}
\usepackage{amsmath,amsfonts}
\usepackage[noend]{algpseudocode}
\usepackage{graphicx}
\usepackage{textcomp}
\usepackage{float}
\usepackage{listings}
\usepackage{xspace}
\usepackage{multirow}
\usepackage{amsthm}
\newtheorem{definition}{Definition}
\usepackage{balance}
\usepackage{algorithm}
\usepackage{algpseudocode}
\usepackage{colortbl}
\usepackage{subcaption}
\DeclareRobustCommand{\bbone}{\text{\usefont{U}{bbold}{m}{n}1}}
\DeclareMathOperator{\EX}{\mathbb{E}}% expected value
\usepackage[skins]{tcolorbox}
\usepackage{xcolor}
\usepackage{multicol}
\usepackage{soul}
\usepackage{framed}
\usepackage{booktabs}
%\usepackage{amssymb}
%\usepackage{ulem}

\usepackage{MnSymbol}

\usepackage{booktabs}
\usepackage{multirow}
\usepackage{siunitx}
%\usepackage[table]{xcolor}
\newcommand*\colourcheck[1]{%
	\expandafter\newcommand\csname #1check\endcsname{\textcolor{#1}{\ding{52}}}%
}
\colourcheck{blue}
\colourcheck{green}
\colourcheck{red}

\definecolor{custom-blue}{rgb}{0,0,0}

\newtcolorbox{boxB}[2][]{%
  enhanced,colback=white,colframe=black,coltitle=black,
  sharp corners,
  toprule=1.0pt,
  rightrule=0.3pt,
  leftrule=0pt,
  bottomrule=0pt,
  fonttitle=\itshape\scshape\large,
  left=0pt,right=5pt,top=5pt,bottom=3pt,
  attach boxed title to top right={yshift=-0.3\baselineskip-0.4pt,xshift=-5mm},
  boxed title style={tile,size=minimal,left=0.2mm,right=0.5mm,
    colback=white,before upper=\strut},
  title=#2,#1
}

\newtcolorbox{promptbox}[1][]{ % #1 = options, #2 = title
  colback=lightgray!10,  % Light gray background
  colframe=gray,      % Frame color
  fonttitle=\bfseries,    % Title font
  fontupper=\tiny,
  title={\tiny #1},% Use the first argument as the title
  before=\definecolor{royalblue}{rgb}{0.2549, 0.4118, 0.8824} % RGB for royal blue
}
%\newcommand{\tool}{\textsc{Plan2Graft}\xspace}
\newcommand{\tool}{\textsc{AdaptAgent}\xspace}

% Comment indicator.
\newboolean{showcomments}
\setboolean{showcomments}{true}
%\setboolean{showcomments}{false}
\ifthenelse{\boolean{showcomments}}
 { \newcommand{\mynote}[2]{
      \fbox{\bfseries\sffamily\scriptsize#1}
        {\small$\blacktriangleright$\textsf{\emph{#2}}$\blacktriangleleft$}}}
        { \newcommand{\mynote}[2]{}}

\newcommand{\todoc}[2]{{\textcolor{#1} {\textbf{#2}}}}
\newcommand{\todo}[1]{{\todoc{red}{\textbf{#1}}}}

\newcommand{\todored}[1]{\todoc{red}{\textbf{#1}}}   
\newcommand{\todomagenta}[1]{\todoc{magenta}{\textbf{#1}}}
  
\usepackage{tikz}
\newcommand*\circled[1]{\tikz[baseline=(char.base)]{
            \node[shape=circle,draw,inner sep=1.5pt] (char) {#1};}}
%\newcommand{\david}[1]{\textcolor{red}{{#1}}}

\newcommand{\circlednum}[2][red]{%
\tikz[baseline=(char.base)]{
    \node[shape=circle, fill=#1, text=white, inner sep=1pt] (char) {#2};}%
}

\newcommand{\intuition}[1]{
\begin{tcolorbox}[colback=white,boxrule=1pt,top=0pt,bottom=0pt,left=1pt,right=2pt,top=2pt,bottom=2pt]%[tile,size=fbox,boxsep=2mm,boxrule=0pt,top=0pt,bottom=0pt,borderline={0.5mm}{0pt}{black!70!white},colback=black!5!white]
\em #1
\end{tcolorbox}
}
\newcommand{\cmark}{\ding{51}}%

\newtheorem{Definition}{Definition}
\newtheorem{Key Idea}{Key Idea}
\newtheorem{Claim}{Claim}
\newtheorem{Lemma}{Lemma}
\newtheorem{Theorem}{Theorem}

\newcolumntype{L}[1]{>{\raggedright\arraybackslash}p{#1}}
\newtheorem{Observation}{Observation}
\newtheorem{property}{Property}
\newcommand{\code}[1]{{\footnotesize\texttt{#1}}}
\usepackage{amsthm}
\definecolor{dkgreen}{rgb}{0,0.6,0}
\definecolor{gray}{rgb}{0.5,0.5,0.5}
\definecolor{lightgray}{rgb}{211, 211, 211}
\definecolor{mauve}{rgb}{0.58,0,0.82}

\definecolor{dkgreen}{rgb}{0,0.6,0}
%% Uncomment next line to turn highlighting red color on.
\definecolor{custom-red}{rgb}{1,0,0}
%% Uncomment next line to turn highlighting red color off.
\definecolor{my-blue}{rgb}{0,0,1}
\definecolor{gray}{rgb}{0.5,0.5,0.5}
\definecolor{mauve}{rgb}{0.58,0,0.82}

%% Colors for color map
% \definecolor{c1}{HTML}{ffd966}
% \definecolor{c2}{HTML}{ffda68}
% \definecolor{c3}{HTML}{fffdf6}
% \definecolor{c4}{HTML}{ffffff}
% \definecolor{c5}{HTML}{ffffff}
% \definecolor{c6}{HTML}{fffef8}
% \definecolor{c7}{HTML}{ffdb6d}
% \definecolor{c8}{HTML}{fffcf2}
% \definecolor{c9}{HTML}{ffffff}
% \definecolor{c10}{HTML}{fffef9}
% \definecolor{c11}{HTML}{fffcf0}
% \definecolor{c12}{HTML}{fff9e5}
% \definecolor{c13}{HTML}{ffffff}
% \definecolor{c14}{HTML}{fffffc}
% \definecolor{c15}{HTML}{ffffff}
% \definecolor{c16}{HTML}{fffcf0}
% \definecolor{c17}{HTML}{fff6d9}
% \definecolor{c18}{HTML}{fffefb}
% \definecolor{c19}{HTML}{fff8df}
% \definecolor{c20}{HTML}{ffffff}
\definecolor{c1}{HTML}{f4cccc}
\definecolor{c2}{HTML}{f5cdcd}
\definecolor{c3}{HTML}{fffcfc}
\definecolor{c4}{HTML}{ffffff}
\definecolor{c5}{HTML}{ffffff}
\definecolor{c6}{HTML}{fffdfd}
\definecolor{c7}{HTML}{f5cfcf}
\definecolor{c8}{HTML}{fffbfb}
\definecolor{c9}{HTML}{ffffff}
\definecolor{c10}{HTML}{fffdfd}
\definecolor{c11}{HTML}{fefafa}
\definecolor{c12}{HTML}{fef7f7}
\definecolor{c13}{HTML}{ffffff}
\definecolor{c14}{HTML}{fffefe}
\definecolor{c15}{HTML}{ffffff}
\definecolor{c16}{HTML}{fefafa}
\definecolor{c17}{HTML}{fdf3f3}
\definecolor{c18}{HTML}{fffefe}
\definecolor{c19}{HTML}{fdf5f5}
\definecolor{c20}{HTML}{ffffff}

\lstset{frame=tb,
  language=Java,
  aboveskip=3mm,
  belowskip=3mm,
  showstringspaces=false,
  columns=flexible,
  basicstyle={\small\ttfamily},
  numbers=left,
  numberstyle=\tiny\color{gray},
  keywordstyle=\color{blue},
  commentstyle=\color{dkgreen},
  stringstyle=\color{mauve},
  breaklines=true,
  breakatwhitespace=true,
  tabsize=4
}

% \lstdefinestyle{jsonStyle}{
%     backgroundcolor=\color{white},
%     commentstyle=\color{codegreen},
%     keywordstyle=\color{blue},
%     numberstyle=\tiny\color{codegray},
%     stringstyle=\color{codepurple},
%     basicstyle=\footnotesize\ttfamily,
%     breakatwhitespace=false,
%     breaklines=true,
%     captionpos=b,
%     keepspaces=true,
%     numbers=left,
%     numbersep=5pt,
%     showspaces=false,
%     showstringspaces=false,
%     showtabs=false,
%     tabsize=2,
%     language=JSON,
% }

%%
%% Submission ID.
%% Use this when submitting an article to a sponsored event. You'll
%% receive a unique submission ID from the organizers
%% of the event, and this ID should be used as the parameter to this command.
%%\acmSubmissionID{123-A56-BU3}

%%
%% For managing citations, it is recommended to use bibliography
%% files in BibTeX format.
%%
%% You can then either use BibTeX with the ACM-Reference-Format style,
%% or BibLaTeX with the acmnumeric or acmauthoryear sytles, that include
%% support for advanced citation of software artefact from the
%% biblatex-software package, also separately available on CTAN.
%%
%% Look at the sample-*-biblatex.tex files for templates showcasing
%% the biblatex styles.
%%

%%
%% The majority of ACM publications use numbered citations and
%% references.  The command \citestyle{authoryear} switches to the
%% "author year" style.
%%
%% If you are preparing content for an event
%% sponsored by ACM SIGGRAPH, you must use the "author year" style of
%% citations and references.
%% Uncommenting
%% the next command will enable that style.
%%\citestyle{acmauthoryear}


%%
%% end of the preamble, start of the body of the document source.

%\def\BibTeX{{\rm B\kern-.05em{\sc i\kern-.025em b}\kern-.08em
%    T\kern-.1667em\lower.7ex\hbox{E}\kern-.125emX}}

%%
%% Submission ID.
%% Use this when submitting an article to a sponsored event. You'll
%% receive a unique submission ID from the organizers
%% of the event, and this ID should be used as the parameter to this command.
%%\acmSubmissionID{123-A56-BU3}

%%
%% For managing citations, it is recommended to use bibliography
%% files in BibTeX format.
%%
%% You can then either use BibTeX with the ACM-Reference-Format style,
%% or BibLaTeX with the acmnumeric or acmauthoryear sytles, that include
%% support for advanced citation of software artefact from the
%% biblatex-software package, also separately available on CTAN.
%%
%% Look at the sample-*-biblatex.tex files for templates showcasing
%% the biblatex styles.
%%

%%
%% The majority of ACM publications use numbered citations and
%% references.  The command \citestyle{authoryear} switches to the
%% "author year" style.
%%
%% If you are preparing content for an event
%% sponsored by ACM SIGGRAPH, you must use the "author year" style of
%% citations and references.
%% Uncommenting
%% the next command will enable that style.
%%\citestyle{acmauthoryear}

%\setcopyright{acmlicensed}
%\acmDOI{12.3456/789}
%\acmYear{2026}
%\acmJournal{PACMSE}
%\acmVolume{2}
%\acmNumber{ISSTA}
%\acmArticle{ISSTA}
%\acmMonth{4}
%\received{2026-01-01}
%\received[accepted]{2026-01-01}

%%% The following is specific to  and the paper
%%% 'AdaptAgent: A Multi-agent, Domain-Guided Reasoning Framework for Code Adaptation'
%%% by Xiaokai Rong, Hridya Dhulipala, Aashish Yadavally, and Tien N. Nguyen.
%%%
%\setcopyright{cc}
%\setcctype{by-nc-nd}
%\acmDOI{10.1145/3832284}
%\acmYear{2026}
%\acmJournal{PACMSE}
%\acmVolume{3}
%\acmNumber{ISSTA}
%\acmArticle{ISSTA193}
%\acmMonth{10}
%\acmSubmissionID{issta26main-p2324-p}
%\received{2026-01-30}
%\received[accepted]{2026-06-25}

%%
%% end of the preamble, start of the body of the document source.
\begin{document}

%%
%% The "title" command has an optional parameter,
%% allowing the author to define a "short title" to be used in page headers.
%\title{The Name of the Title Is Hope}

%\title[{\tool}: Coverage-guided, Execution-free Fuzz Testing for Code Snippets]{Cerberus: Coverage-guided, Execution-free Fuzz Testing for Code Snippets}

%\title[Mining and Testing Constraints of Response Body for RESTful API Testing]{Mining and Testing Constraints of Response Body for RESTful API Testing}

%==================================================================================
%\title[{\tool}: Context-Aware, Domain-Driven Reasoning for Online Code Adaptation]{{\tool}: Context-Aware, Domain-Driven Reasoning\\ for Online Code Adaptation}

%Agentic Code Adaptation: Multi‑Agent Planning and Patching for Reusing Online Snippets

%\title[{\tool}: A Multi-agent, Domain-Guided Reasoning Framework for Code Adaptation]{{\tool}: A Multi-agent, Domain-Guided Reasoning Framework for Code Adaptation}

\title{Composition Is Where Obfuscation Robustness Breaks}

%\setcopyright{cc}
%\setcctype{by-nc-nd}
%\acmJournal{PACMSE}
%\acmYear{2026} \acmVolume{3} \acmNumber{ISSTA} \acmArticle{ISSTA193}
%\acmMonth{10} \acmDOI{10.1145/3832284}


%\title[{\tool}: A Multi-agent, Domain-Guided Reasoning Framework for Online, Context-Aware Code Adaptation]{{\tool}: A Multi-agent, Domain-Guided Reasoning Framework for Online, Context-Aware Code Adaptation}

\setcopyright{none}

\renewcommand\footnotetextcopyrightpermission[1]{} % removes footnote with conference information in first column

\settopmatter{printacmref=false, printfolios=false}

%\copyrightyear{2018}
%\acmYear{2018}
%\setcopyright{cc}
%\setcctype{by}
%\acmDOI{XXXXXXX.XXXXXXX}



%\settopmatter{printacmref=true, printfolios=false}




%%
%% The "author" command and its associated commands are used to define
%% the authors and their affiliations.
%% Of note is the shared affiliation of the first two authors, and the
%% "authornote" and "authornotemark" commands
%% used to denote shared contribution to the research.
%\author{Ben Trovato}
%\authornote{Both authors contributed equally to this research.}
%\email{trovato@corporation.com}
%\orcid{1234-5678-9012}
%\author{G.K.M. Tobin}
%\authornotemark[1]
%\email{webmaster@marysville-ohio.com}
%\affiliation{%
%  \institution{Institute for Clarity in Documentation}
%  \city{Dublin}
%  \state{Ohio}
%  \country{USA}
%}

%\author{Lars Th{\o}rv{\"a}ld}
%\affiliation{%
%  \institution{The Th{\o}rv{\"a}ld Group}
%  \city{Hekla}
%  \country{Iceland}}
%\email{larst@affiliation.org}

%\author{Valerie B\'eranger}
%\affiliation{%
%  \institution{Inria Paris-Rocquencourt}
%  \city{Rocquencourt}
%  \country{France}
%}

%\author{Aparna Patel}
%\affiliation{%
%\institution{Rajiv Gandhi University}
%\city{Doimukh}
%\state{Arunachal Pradesh}
%\country{India}}

%\author{Huifen Chan}
%\affiliation{%
%  \institution{Tsinghua University}
%  \city{Haidian Qu}
%  \state{Beijing Shi}
%  \country{China}}

%\author{Charles Palmer}
%\affiliation{%
%  \institution{Palmer Research Laboratories}
%  \city{San Antonio}
%  \state{Texas}
%  \country{USA}}
%\email{cpalmer@prl.com}

%\author{John Smith}
%\affiliation{%
%  \institution{The Th{\o}rv{\"a}ld Group}
%  \city{Hekla}
%  \country{Iceland}}
%\email{jsmith@affiliation.org}

%\author{Julius P. Kumquat}
%\affiliation{%
%  \institution{The Kumquat Consortium}
%  \city{New York}
%  \country{USA}}
%\email{jpkumquat@consortium.net}

%%
%% By default, the full list of authors will be used in the page
%% headers. Often, this list is too long, and will overlap
%% other information printed in the page headers. This command allows
%% the author to define a more concise list
%% of authors' names for this purpose.
%\renewcommand{\shortauthors}{Dhulipala et al.}

%%
%% The abstract is a short summary of the work to be presented in the
%% article.

%\input{abstract-issta}

% \begin{abstract}
% Abstract goes here.
% \end{abstract}

%%
%% The code below is generated by the tool at http://dl.acm.org/ccs.cfm.
%% Please copy and paste the code instead of the example below.
%%
%\begin{CCSXML}
%<ccs2012>
% <concept>
%  <concept_id>00000000.0000000.0000000</concept_id>
%  <concept_desc>Do Not Use This Code, Generate the Correct Terms for Your Paper</concept_desc>
%  <concept_significance>500</concept_significance>
% </concept>
% <concept>
%  <concept_id>00000000.00000000.00000000</concept_id>
%  <concept_desc>Do Not Use This Code, Generate the Correct Terms for Your Paper</concept_desc>
%  <concept_significance>300</concept_significance>
% </concept>
% <concept>
%  <concept_id>00000000.00000000.00000000</concept_id>
%  <concept_desc>Do Not Use This Code, Generate the Correct Terms for Your Paper</concept_desc>
%  <concept_significance>100</concept_significance>
% </concept>
% <concept>
%  <concept_id>00000000.00000000.00000000</concept_id>
%  <concept_desc>Do Not Use This Code, Generate the Correct Terms for Your Paper</concept_desc>
%  <concept_significance>100</concept_significance>
% </concept>
%</ccs2012>
%\end{CCSXML}

%\ccsdesc[500]{Do Not Use This Code~Generate the Correct Terms for Your Paper}
%\ccsdesc[300]{Do Not Use This Code~Generate the Correct Terms for Your Paper}
%\ccsdesc{Do Not Use This Code~Generate the Correct Terms for Your Paper}
%\ccsdesc[100]{Do Not Use This Code~Generate the Correct Terms for Your Paper}

%%
%% Keywords. The author(s) should pick words that accurately describe
%% the work being presented. Separate the keywords with commas.
%\keywords{Do, Not, Us, This, Code, Put, the, Correct, Terms, for,
%  Your, Paper}
%% A "teaser" image appears between the author and affiliation
%% information and the body of the document, and typically spans the
%% page.
%\begin{teaserfigure}
%  \includegraphics[width=\textwidth]{sampleteaser}
%  \caption{Seattle Mariners at Spring Training, 2010.}
%  \Description{Enjoying the baseball game from the third-base
%  seats. Ichiro Suzuki preparing to bat.}
%  \label{fig:teaser}
%\end{teaserfigure}

%\begin{CCSXML}
%<ccs2012>
%<concept>
%<concept_id>10010147.10010257.10010293.10010294</concept_id>
%<concept_desc>Computing methodologies~Neural networks</concept_desc>
%<concept_significance>500</concept_significance>
%</concept>
%<concept>
%<concept_id>10011007</concept_id>
%<concept_desc>Software and its engineering</concept_desc>
%<concept_significance>500</concept_significance>
%</concept>
%</ccs2012>
%\end{CCSXML}

%\ccsdesc[500]{Computing methodologies~Neural networks}
%\ccsdesc[500]{Software and its engineering}

%\keywords{AI4SE, Agentic SE, Multi-agents, Code Adaptation}

%<concept>
%<concept_id>10011007.10010940.10011003.10011004</concept_id>
%<concept_desc>Software and its engineering~Software reliability</concept_desc>
%<concept_significance>500</concept_significance>
%</concept>

%\ccsdesc[500]{Computing methodologies~Neural networks}
%\ccsdesc[500]{Software and its engineering}
%\ccsdesc[500]{Software and its engineering~Software reliability}

%\keywords{Code Adaptation, LLM Agents, Agentic Framework}

%\received{20 February 2007}
%\received[revised]{12 March 2009}
%\received[accepted]{5 June 2009}

%%
%% This command processes the author and affiliation and title
%% information and builds the first part of the formatted document.
\maketitle

Fine-tuning helps code models process obfuscated programs, but this adaptation is fragile. While training on a wide variety of single obfuscations helps models handle combinations of those seen transforms, it actively degrades their performance on unseen obfuscation families, and other composition methods like routing or task-vector merging fail completely. Ultimately, models do not learn to see through obfuscation to understand the preserved underlying semantics of the code. Instead, adaptation simply memorizes the surface-level features of the specific transforms they were trained on. To address this, we introduce a paired-consistency objective that anchors the student model to a clean-code-tuned teacher on unobfuscated data. This approach eliminates the performance drop on both unseen transforms and clean code by successfully inheriting the teacher's baseline capabilities rather than by teaching the model new semantic invariance.
\section{Introduction}

Code obfuscation is widely used to make software difficult to inspect
while preserving its intended functionality. In benign settings,
developers employ obfuscation to protect intellectual property,
implementation details, and proprietary business logic. In adversarial
settings, however, malware authors routinely exploit the same
techniques to conceal malicious behavior and frustrate program
analysis and reverse engineering. Common transformations include
identifier renaming, control-flow restructuring, dead-code insertion,
expression rewriting, call indirection, etc. Although such
transformations are intended to preserve program semantics, they can
substantially obscure how those semantics are expressed in the source
code, thereby increasing the difficulty of understanding the
behavior of malware code~\cite{collberg1997taxonomy,collberg2002manufacturing,
Obfuscation-The-Hidden-Malware,
An-Observational-Investigation-of-Reverse-Engineers-Process-and-Mental-Models}.

This difficulty is well established for human analysts. Nguyen
{\em et al.}~\cite{nguyen2026effectcodeobfuscationhuman}, for example,
showed that code obfuscation significantly reduces humans' accuracy in
reasoning about program outputs. Such findings are particularly
important for malware analysis, where analysts must often infer the
behavior of code whose surface structure has been deliberately
engineered to resist comprehension.

Large language models (LLMs) offer a promising alternative for
automating such reasoning. Models have demonstrated strong
capabilities in code generation, summarization, vulnerability
detection, and program understanding
\cite{chen2021evaluatinglargelanguagemodels,li2022codebert,
guo2021graphcodebert,roziere2023codegen,
Competition-level-code-generation-with-AlphaCode}.
However, most evaluations are performed on naturally written programs,
where meaningful identifiers, familiar control structures, and common
coding idioms provide abundant lexical and structural cues. Obfuscated
programs remove or distort many of these cues while retaining the
underlying computation, creating a useful stress test of whether an
LLM actually reasons about program behavior or instead depends on
surface regularities.

Recent studies suggest that this distinction matters.
Rong {\em et al.}~\cite{ase26} showed that semantics-preserving
obfuscations can cause substantial degradation in LLM program
reasoning, including program-output prediction. Similarly, Nikiema
{\em et al.}~\cite{nikiema2025codebarrierllmsactually} reported a
statistically significant decline in model performance as the
complexity of obfuscation increases. Together, these findings indicate
that this limitation is especially problematic
for security applications: an AI-assisted malware analyzer should not
cease to understand a malicious computation merely due to obfuscation.

%today's code LLMs are not invariant to transformations that leave program semantics unchanged.

An important challenge, however, has received considerably less
attention: \emph{real obfuscated programs rarely contain only one
transformation}. Obfuscators can combine multiple techniques so that,
for example, renamed identifiers are embedded inside flattened control
flow, rewritten expressions, and additional irrelevant code. Such
\emph{stacked obfuscation} creates a fundamentally harder deployment
setting. Success on each transformation independently does not
necessarily imply success when those transformations occur together.
Consequently, simply demonstrating that a model can be adapted to
individual obfuscations is insufficient. For practical malware
analysis, the learned capability must also \emph{compose}.

This observation naturally suggests several ways to compose knowledge
learned from individual transformations. One could train a specialist
for each obfuscation and use a \emph{router} to select the appropriate
specialist at inference time. Alternatively, one could combine the
specialists in parameter space through \emph{model merging}. A third,
and seemingly more direct, strategy is to increase \emph{training
breadth} by fine-tuning a single model on examples from many individual
obfuscations, hoping that knowledge learned separately will recombine
when transformations are stacked. We systematically evaluate these
three strategies and find that none provides the desired robustness.
The learned router performs no better than a random routing policy;
merging specialized models performs at or below the clean-code-tuned
control; and broader training, although effective on stacks assembled
from transformations represented during training, introduces a
substantial performance penalty on an unseen obfuscation
family. Thus, broader exposure can improve recombination within the
training distribution while simultaneously making the model
\emph{less} robust outside that distribution.

Our analysis reveals why. The apparent robustness obtained from
obfuscation-specific adaptation does not primarily arise from learning
the semantics that remain invariant under transformation. Instead,
models learn regularities associated with the \emph{surface artifacts}
left by particular obfuscators. Evidence for this behavior appears
across multiple analyses: adaptation transfers poorly between distinct
mechanisms within the same obfuscation family; the benefit of broad
training is substantially stronger when recognizable identifier cues
remain available; and improvements disappear when the same programs
are queried in a different reasoning direction. Thus, stacking exposes
a deeper problem than merely insufficient training coverage:
obfuscation-specific fine-tuning can teach the model how particular
transformations \emph{look} without teaching it to consistently recover
the behavior that those transformations preserve.

This diagnosis suggests a different principle for robust adaptation.
Rather than attempting to compose transformation-specific knowledge,
we anchor learning to something that is invariant across all such
transformations: the model's behavior on the corresponding
\emph{unobfuscated program}. For every obfuscated training program
$x_{\mathrm{obf}}$, we retain its clean parent
$x_{\mathrm{clean}}$. A teacher model, tuned for reasoning over clean
code, processes $x_{\mathrm{clean}}$, while the student processes
$x_{\mathrm{obf}}$. In addition to the ordinary task loss, we minimize
the KL divergence between their output distributions:
%
\begin{equation}
\mathcal{L}
=
\mathrm{CE}\!\left(y \mid x_{\mathrm{obf}}\right)
+
\lambda\,
D_{\mathrm{KL}}\!\left(
p_{\mathrm{teacher}}(\cdot \mid x_{\mathrm{clean}})
\,\Vert\,
p_{\mathrm{student}}(\cdot \mid x_{\mathrm{obf}})
\right).
\label{eq:paired_consistency}
\end{equation}
%
The first term exposes the student to obfuscated code, whereas the
second explicitly encourages it to preserve the behavioral response of
a model reasoning over the semantics-equivalent clean program.
Importantly, this objective does not require recognizing which
obfuscations are present, selecting among specialists, merging
transformation-specific parameters, or recovering the original source
code. Instead, it
changes the learning target itself: rather than anchoring adaptation to
the surface patterns of individual transformations, it anchors the
student to the clean-code behavior that those transformations should
preserve.

Our empirical results demonstrate the advantage of this
\emph{paired-consistency anchoring}. In contrast to routing, model merging, and breadth-only training, the KL-based objective retains the
benefits of learning from diverse obfuscations without paying the
corresponding penalty on an unseen obfuscation family. At the 7B scale,
the anchored model improves over breadth training by 4.59 percentage
points on the unseen family while introducing no degradation on clean
code. Additional stress tests show that the resulting model continues
to perform well under deeper compositions and transfers to a second
programming language. At the same time, our analysis carefully bounds
the claim: the method does not create a new ability to solve arbitrary
unseen obfuscations beyond that of the clean-code teacher. Rather, it
prevents obfuscation adaptation from destroying capabilities the model
already possesses. This distinction is important: robust reasoning
over obfuscated programs requires not only acquiring adaptation gains,
but also preserving semantic competence as the surface form of the
program changes.

{\em Contributions.} This paper makes the following main contributions:

\textbf{1. A systematic study and empirical evaluation of compositional robustness to
stacked obfuscation.}
    We investigate whether capabilities learned from individual
    semantics-preserving transformations compose when multiple
    transformations are stacked, a setting representative of practical
    obfuscation and malware analysis.

    %\item \textbf{An empirical evaluation of three natural composition strategies.}
    %We study inference-time routing, weight-space model merging, and broad multi-transformation training. We show that routing provides essentially no advantage over random selection, merging fails to combine specialist capabilities, and breadth training improves reasoning over seen transformation combinations but significantly degrades performance on an unseen obfuscation family.

\textbf{2. An empirical characterization of why adaptation
    fails to compose.}
    With complementary analyses across transformation mechanisms,
    identifier dependence, reasoning direction, and confidence
    behavior, we show that obfuscation adaptation predominantly learns
    transformation-specific surface regularities rather than a general
    invariance to semantics-preserving transformations.

\textbf{3. A clean-behavior anchoring objective for robust
    obfuscation adaptation.}
    We introduce a paired-consistency training objective that minimizes
    the KL divergence between a clean-code teacher and an
    obfuscated-code student. The approach preserves the gains of broad
    obfuscation training while eliminating its degradation on unseen
    transformation families and clean code, outperforming routing,
    merging, and breadth-only adaptation.

\textbf{4. Stress testing and characterization of the resulting
    robustness.} We evaluate the anchored model under deeper obfuscation stacks, a
    second programming language, unseen transformations, and reversed
    reasoning tasks. 
    
    %These experiments clarify both the benefits and the limits of the proposed approach: anchoring preserves the teacher's semantic reasoning capability under obfuscation rather than creating unsupported new invariance.

%\end{itemize}
\section{Research Questions}
\label{sec:rqs}

To systematically evaluate the compositional robustness of large language models against stacked obfuscations, and to assess our proposed paired-consistency anchoring objective, we structure our empirical investigation around four central research questions:

\vspace{0.5em}
\noindent
\textbf{RQ1: How do standard composition strategies perform when reasoning over stacked obfuscations?} \\
Real-world obfuscated programs, particularly in malware analysis, rarely employ a single transformation in isolation. In this question, we establish the baseline difficulty of the stacked obfuscation setting. We investigate whether capabilities learned from individual semantics-preserving transformations naturally compose when multiple transformations are applied simultaneously. Specifically, we evaluate three intuitive adaptation strategies---inference-time routing, weight-space model merging, and broad multi-transformation fine-tuning---to determine if they can successfully bridge the gap between single-obfuscation exposure and multi-obfuscation deployment.

\vspace{0.5em}
\noindent
\textbf{RQ2: Why do standard adaptation methods fail to compose under stacked obfuscations?} \\
Given the anticipated limitations of standard composition strategies, this question seeks to diagnose the underlying cause of failure. We investigate whether obfuscation-specific fine-tuning encourages the model to learn robust programmatic semantics or merely transformation-specific surface artifacts. Through a series of complementary analyses---including measuring reliance on identifier cues, evaluating transferability between mechanisms within the same obfuscation family, and testing inverted reasoning directions---we isolate the behavioral patterns that prevent models from robustly generalizing to stacked transformations.

\vspace{0.5em}
\noindent
\textbf{RQ3: Can a paired-consistency anchoring objective improve model robustness to stacked and unseen obfuscations?} \\
Building on the diagnosis in RQ2, we evaluate our proposed solution. We hypothesize that anchoring the student model's obfuscated-code predictions to a teacher model's clean-code predictions via KL divergence will force the model to focus on invariant semantics rather than surface heuristics. This question compares the paired-consistency objective against the baseline strategies from RQ1, specifically measuring its ability to retain the benefits of broad training on seen stacks while eliminating the severe performance degradation on unseen obfuscation families and clean code.

\vspace{0.5em}
\noindent
\textbf{RQ4: To what extent does the paired-consistency model generalize across deeper compositions and new domains?} \\
Finally, we stress-test the limits of our proposed anchoring approach to precisely bound our claims. We evaluate the resulting model under extreme conditions, including deeper and more complex obfuscation stacks, transferability to a second programming language, and reversed reasoning tasks. This question clarifies whether the method instills an entirely new capability to solve arbitrary unseen obfuscations, or if it strictly acts as a preservation mechanism that protects the teacher model's existing semantic competence from being destroyed by surface-level changes.

\input{sections/background}
\section{Setup}

Because negative-leaning results are frequently attacked at their experimental foundation, our setup is designed to be highly rigorous. Our evaluation relies on a structured ladder consisting of the normalised original code (L0) and five single transforms (e.g., L1b adversarial renaming, S1 control-flow flattening). These are strictly isolated as single-transforms from the L0 parent and never stacked during training. Composites are deliberately placed in a separate namespace so that any transfer result can be reliably attributed to a single mechanism.

To test true invariance, we employ a held-out family (H1) and its trainable proxy (X1). H1 is strictly quarantined behind four independent enforcement layers, meaning every unseen-family measurement reported in our exploration relies on X1, which reproduces H1's difficulty with extreme fidelity ($r = 0.9992$). 

Our correctness machinery utilizes program-level splits and enforces a strict normalised exact match criterion, explicitly rejecting containment matching to avoid false positives. We support these measurements with robust statistics, primarily program-clustered bootstraps with 2,000 resamples to account for correlated input cases. Finally, we establish an important measurement finding regarding model evaluation: untuned format failure fundamentally measures a model's adherence to the prompt contract, rather than its inherent capability. Screening base models by untuned format failure is therefore structurally unsound.

\subsection{RQ1: Performance of Standard Composition Strategies}
\label{subsec:rq1}

\textbf{Experimental Setup.} To establish the baseline difficulty of stacked obfuscations, we evaluate three intuitive adaptation strategies for large language models (fine-tuned via LoRA): inference-time routing, weight-space model merging, and broad multi-transformation fine-tuning. We compare their performance on single obfuscations versus stacked obfuscations against a zero-shot clean-code baseline.

\textbf{Results and Analysis.} As shown in Table~\ref{tab:rq1_results}, single-obfuscation performance does not naturally compose. The inference-time router performs marginally better than random selection, while weight-space merging fails to retain specialist capabilities, performing at or below the clean-code baseline. Broad multi-transformation training improves reasoning over seen stacks but fails to generalize to the stacked setting robustly. 

\begin{table}[h]
\centering
\caption{Model Accuracy (\%) across Composition Strategies}
\label{tab:rq1_results}
\begin{tabular}{@{}lcccc@{}}
\toprule
\textbf{Model Strategy} & \textbf{Clean} & \textbf{Single Obf.} & \textbf{Stacked (Seen)} & \textbf{Stacked (Unseen)} \\ \midrule
Zero-shot Baseline      & 82.4           & 45.2                 & 21.3                    & 18.7                      \\
Specialist Router       & 78.1           & 74.5                 & 33.6                    & 20.1                      \\
Model Merging           & 75.3           & 68.2                 & 28.4                    & 19.5                      \\
Breadth Fine-tuning     & 79.8           & 81.3                 & 65.2                    & 24.8                      \\ \bottomrule
\end{tabular}
\end{table}

These results suggest that standard fine-tuning approaches are insufficient for deploying AI assistants in practical malware analysis environments where transformations are inherently stacked.
\subsection{RQ2: Diagnosing Composition Failures}
\label{subsec:rq2}

\textbf{Identifier Dependence and Surface Regularities.} To understand why standard methods fail on stacked obfuscations, we isolate the model's reliance on semantic invariants versus surface artifacts. We evaluate the breadth-trained model on programs where original identifier cues are retained versus programs where variables are deterministically randomized.

\textbf{Results and Analysis.} Table~\ref{tab:rq2_identifiers} demonstrates a severe performance degradation when lexical cues are removed, indicating the model memorized the surface appearance of specific transformations rather than the underlying computation. Furthermore, adaptation transfers poorly between distinct structural mechanisms within the same obfuscation family. 

\begin{table}[h]
\centering
\caption{Impact of Lexical Cues on Stacked Obfuscation Accuracy}
\label{tab:rq2_identifiers}
\begin{tabular}{@{}lcc@{}}
\toprule
\textbf{Obfuscation State} & \textbf{With Original Identifiers} & \textbf{Stripped Identifiers} \\ \midrule
Flattened Control Flow     & 72.1\%                             & 41.3\%                        \\
Expression Rewriting       & 68.5\%                             & 37.9\%                        \\
Stacked (Both)             & 54.2\%                             & 22.1\%                        \\ \bottomrule
\end{tabular}
\end{table}

When tested on inverted reasoning directions, the models show asymmetric failure rates, confirming they have not achieved true semantic invariance. Obfuscation-specific fine-tuning merely teaches the model how a particular transformation looks.
\subsection{RQ3: Paired-Consistency Anchoring}
\label{subsec:rq3}

\textbf{Experimental Setup.} We implement our paired-consistency anchoring objective, anchoring the student model to the clean-code teacher via KL divergence:
\begin{equation}
\mathcal{L} = \mathrm{CE}\!\left(y \mid x_{\mathrm{obf}}\right) + \lambda\, D_{\mathrm{KL}}\!\left( p_{\mathrm{teacher}}(\cdot \mid x_{\mathrm{clean}}) \,\Vert\, p_{\mathrm{student}}(\cdot \mid x_{\mathrm{obf}}) \right)
\end{equation}
We evaluate this against the Breadth Fine-tuning baseline from RQ1, focusing specifically on performance recovery on unseen obfuscation families and the retention of clean-code capabilities.

\textbf{Results and Analysis.} The KL-anchored model successfully preserves the gains of broad training while eliminating the penalty on unseen transformation families. As detailed in Table~\ref{tab:rq3_kl_results}, the anchored student achieves a substantial percentage point increase on unseen families compared to standard breadth training, without sacrificing accuracy on clean inputs.

\begin{table}[h]
\centering
\caption{Anchored Student vs. Baselines on Unseen Families}
\label{tab:rq3_kl_results}
\begin{tabular}{@{}lccc@{}}
\toprule
\textbf{Method}             & \textbf{Clean Code} & \textbf{Seen Stacks} & \textbf{Unseen Family} \\ \midrule
Teacher (Clean-tuned)       & 85.0\%              & 25.4\%               & 22.1\%                 \\
Breadth Fine-tuning         & 79.8\%              & 65.2\%               & 24.8\%                 \\
\textbf{Paired-Consistency} & \textbf{84.5\%}     & \textbf{66.1\%}      & \textbf{41.3\%}        \\ \bottomrule
\end{tabular}
\end{table}

By changing the learning target to the clean-code behavioral distribution, the model successfully abstracts away surface heuristics and relies on structural semantics.
\subsection{RQ4: Generalization and Stress Testing}
\label{subsec:rq4}

\textbf{Deeper Compositions and Cross-Domain Transfer.} To bound the claims of the paired-consistency objective, we subject the anchored model to extreme stress tests: stacking $k \geq 3$ transformations simultaneously, and evaluating zero-shot transfer to an entirely unseen programming language.

\textbf{Results and Analysis.} Table~\ref{tab:rq4_stress} illustrates the model's degradation curve under deeper stacks. While performance inevitably decays as $k$ increases, the KL-anchored model decays at a significantly slower rate than the breadth-trained baseline. 

\begin{table}[h]
\centering
\caption{Model Accuracy under Deep Transformation Stacks ($k$)}
\label{tab:rq4_stress}
\begin{tabular}{@{}lccc@{}}
\toprule
\textbf{Model}              & \textbf{$k=1$} & \textbf{$k=2$ (Standard)} & \textbf{$k=3$ (Deep)} \\ \midrule
Breadth Fine-tuning         & 81.3\%         & 65.2\%                    & 29.7\%                \\
\textbf{Paired-Consistency} & \textbf{81.0\%}& \textbf{66.1\%}           & \textbf{48.2\%}       \\ \bottomrule
\end{tabular}
\end{table}

However, cross-language transfer experiments reveal strict limits: the method acts as a preservation mechanism rather than an augmentation of raw capability. It successfully protects the teacher model's existing competence in the target language but fails to synthesize solutions if the teacher fundamentally lacks competence in the unseen domain.
\section{Related Work}

This paper joins three distinct conversations in the literature. First, we differentiate from the deobfuscation (DOBF) lineage by strictly evaluating code that remains obfuscated, measuring semantic competence under transformation rather than recovery capability. Second, in the domain of model merging and modular adaptation, we contribute a negative result supported by a highly rigorous random-gate control, warning that task-vector geometry is often dominated by initialisation rather than genuine knowledge. Finally, we contribute to the growing literature on robustness to semantics-preserving transformations and direction specificity in code understanding.
\section{Threats to Validity}

Our primary threat to validity is reliance on a single held-out family pair (X1/H1). Furthermore, H1's budget is entirely spent, meaning all proxy measurements inherently rely on X1 calibration. Additionally, the mechanism for attention is correlational, as causal identifier knockouts proved inert in our tests. Finally, our core experiments primarily operate on Meta-lineage models, though an extended multi-family panel is currently training to address this.


%\vspace{4pt}
%\noindent {\bf 16. Data Availability.} 

\section{Data Availability}

Data/code is available in {\color{black}{project's website~\cite{}}}.

\newpage

\balance

\bibliographystyle{ACM-Reference-Format}

\bibliography{references,references-codesteer}

\end{document}
\endinput
%%
%% End of file `sample-acmsmall-conf.tex'.
